\begin{frame}{Moving Towards Larger Context Windows}
  \begin{itemize}
    \item[\textcolor{red}{$\blacktriangleright$}] \textbf{Efficient Attention Variants:}
      \begin{itemize}
        \item Longformer, BigBird, FlashAttention-2
      \end{itemize}
    \item[\textcolor{red}{$\blacktriangleright$}] \textbf{Memory-Augmented Models:}
      \begin{itemize}
        \item Retrieval-augmented generation (RAG): fetch relevant passages and put them in the prompt (covered later)
        \item Memory layers (e.g., product-key memory)
      \end{itemize}
    \item[\textcolor{red}{$\blacktriangleright$}] \textbf{Chunking + Re-ranking:}
      \begin{itemize}
        \item Process in parts, summarize/join results
      \end{itemize}
    \item[\textcolor{red}{$\blacktriangleright$}] \textbf{Recurrence or State Passing:}
      \begin{itemize}
        \item Transformer-XL, RWKV, RetNet: carry a fixed-size state forward instead of attending to every past token
      \end{itemize}
  \end{itemize}
\end{frame}

\begin{frame}{Moving Towards Larger Context Windows}
  \begin{itemize}
    \item[\textcolor{red}{$\blacktriangleright$}] Serving a long context means keeping the \textbf{KV cache}: the keys and values already computed for every token so far, so that generating the next token does not recompute the whole prompt.
    \item[\textcolor{red}{$\blacktriangleright$}] The cache grows linearly with context length, so memory, not attention cost, is what usually caps the window in practice.
  \end{itemize}
  \vspace{0.5em}
  \begin{center}
  \begin{tikzpicture}[node distance=1.4cm, >=latex, font=\sffamily, scale=0.8, every node/.style={transform shape}]

    % Retriever Node
    % Corner radius must stay below half the node height or the border degenerates.
    \node[draw, fill=purple!20, minimum width=1.5cm, minimum height=0.7cm, rounded corners=0.2cm] (retriever) {Retriever};

    % Long Context LLM container. Its width must exceed the two inner boxes plus the
    % gap between them, or they break out through the sides.
    \node[draw, fill=gray!8, minimum width=6.9cm, minimum height=2.0cm, rounded corners=0.2cm, right=2.0cm of retriever, anchor=west] (llmbox) {};

    % KV Cache box inside LLM
    \node[draw, fill=cyan!15, text width=2.7cm, align=center, minimum height=1.1cm, rounded corners=0.13cm, left=0.1cm of llmbox.center, anchor=east] (kv) {\textbf{KV Cache}\\ 1M{+} tokens};

    % Prompt/User Conversation
    \node[draw, fill=yellow!25, text width=2.7cm, align=center, minimum height=1.1cm, rounded corners=0.13cm, right=0.2cm of kv, anchor=west] (prompt) {\textbf{Prompt}\\ User Conversation};

    % Label above LLM box
    \node[font=\small] at (llmbox.north) [yshift=0.7em] {Long Context LLM};

    % Output node
    \node[draw, fill=gray!15, minimum width=1.0cm, minimum height=0.6cm, rounded corners=0.18cm, right=1.6cm of llmbox, anchor=west] (output) {Output};

    % Connections
    \draw[->, thick] (retriever) -- (llmbox.west);

    \draw[->, thick] (llmbox.east) -- (output);

    % Divider line between KV and Prompt
    \draw[gray!70, thick] ([xshift=0.08cm]kv.east|-kv.north) -- ([xshift=0.08cm]kv.east|-kv.south);

  \end{tikzpicture}
  \end{center}
\end{frame}
